% GENERATED FILE. Edit canonical lessons, then run npm run generate:books.
% canonical-source-hash: fnv1a64:26c436044da68630
% canonical-lessons: UR-C397-fasla, UR-C397-kilomitar, UR-C397-darra, UR-C397-morna, UR-C397-ghumana

\chapter{Distances and Passes}
\label{ch:ur-distances-and-passes}
\hlchaptermodality{397}

\begin{chapteropening}
\textbf{By the end of this chapter:} \emph{I can say how far a place is and which way to turn on a mountain road.}

Complete the last lesson of chapter 397: I can say how far a place is and which way to turn on a mountain road.
\end{chapteropening}

\section[\texorpdfstring{\ur{فاصلہ} (fāslā)}{فاصلہ (fāslā)}]{\ur{فاصلہ} (fāslā) — a distance}

\label{lesson:UR-C397-fasla}



\begin{quote}
Before the new one: say the Urdu for a convoy, then the Urdu for to drive (animals or a cart).
\end{quote}

\subsection*{You'll want to know: \ur{فاصلہ}}
\textbf{\ur{فاصلہ}} — \emph{fāslā} — \textquotedblleft{}a distance\textquotedblright{}.

\begin{quote}
\ur{اسٹیشن} \ur{سے} \ur{ہوٹل} \ur{کا} \ur{فاصلہ} \ur{کم} \ur{ہے۔}
\end{quote}

\begin{quote}
\emph{isṭeshan se hoṭal kā fāslā kam hai.}
\end{quote}

\begin{quote}
It isn't far from the station to the hotel.
\end{quote}

Masculine.

\begin{cousinweb}
\textbf{\ur{فاصلہ}} (\emph{fāslā}) is usually traced, through Persian, to Arabic \emph{fāṣila}, \textquotedblleft{}an interval, a space between\textquotedblright{}.
\end{cousinweb}

\begin{grammarlens}[title={What you've built}]
One more word for getting around.
\end{grammarlens}

\subsection*{Your turn}
\begin{itemize}
\raggedright
  \item \emph{Say it:} \emph{fāslā}
  \item \emph{Say it:} \emph{fāslā}, once more
  \item \emph{Say it:} \emph{isṭeshan se hoṭal kā fāslā kam hai.} — It isn't far from the station to the hotel.
  \item \emph{Recall:} say the Urdu for a convoy, then the Urdu for to drive (animals or a cart), then say \emph{fāslā} again
\end{itemize}

\begin{tcolorbox}[breakable,colback=teal!4,colframe=teal!35!black,title={Before you move on}]
What does \ur{فاصلہ} mean? (\textquotedblleft{}A distance\textquotedblright{}.) Say it once more.
\end{tcolorbox}
\section[\texorpdfstring{\ur{کلومیٹر} (kilomīṭar)}{کلومیٹر (kilomīṭar)}]{\ur{کلومیٹر} (kilomīṭar) — a kilometre}

\label{lesson:UR-C397-kilomitar}



\begin{quote}
Before the new one: say the Urdu for to drive (animals or a cart), then the Urdu for a distance.
\end{quote}

\subsection*{You'll want to know: \ur{کلومیٹر}}
\textbf{\ur{کلومیٹر}} — \emph{kilomīṭar} — \textquotedblleft{}a kilometre\textquotedblright{}.

\begin{quote}
\ur{یہاں} \ur{سے} \ur{شہر} \ur{پچاس} \ur{کلومیٹر} \ur{دور} \ur{ہے۔}
\end{quote}

\begin{quote}
\emph{yahāṅ se shahr pachās kilomīṭar dūr hai.}
\end{quote}

\begin{quote}
The city is fifty kilometres from here.
\end{quote}

Masculine.

\begin{cousinweb}
\textbf{\ur{کلومیٹر}} (\emph{kilomīṭar}) is borrowed from English \emph{kilometre}.
\end{cousinweb}

\begin{grammarlens}[title={What you've built}]
One more word for getting around.
\end{grammarlens}

\subsection*{Your turn}
\begin{itemize}
\raggedright
  \item \emph{Say it:} \emph{kilomīṭar}
  \item \emph{Say it:} \emph{kilomīṭar}, once more
  \item \emph{Say it:} \emph{yahāṅ se shahr pachās kilomīṭar dūr hai.} — The city is fifty kilometres from here.
  \item \emph{Recall:} say the Urdu for to drive (animals or a cart), then the Urdu for a distance, then say \emph{kilomīṭar} again
\end{itemize}

\begin{tcolorbox}[breakable,colback=teal!4,colframe=teal!35!black,title={Before you move on}]
What does \ur{کلومیٹر} mean? (\textquotedblleft{}A kilometre\textquotedblright{}.) Say it once more.
\end{tcolorbox}
\section[\texorpdfstring{\ur{درہ} (darrā)}{درہ (darrā)}]{\ur{درہ} (darrā) — a mountain pass}

\label{lesson:UR-C397-darra}



\begin{quote}
Before the new one: say the Urdu for a distance, then the Urdu for a kilometre.
\end{quote}

\subsection*{You'll want to know: \ur{درہ}}
\textbf{\ur{درہ}} — \emph{darrā} — \textquotedblleft{}a mountain pass\textquotedblright{}.

\begin{quote}
\ur{سردیوں} \ur{میں} \ur{درہ} \ur{برف} \ur{سے} \ur{بند} \ur{ہوتا} \ur{ہے۔}
\end{quote}

\begin{quote}
\emph{sardiyoṅ meṅ darrā barf se band hotā hai.}
\end{quote}

\begin{quote}
In winter the pass is closed by snow.
\end{quote}

Masculine.

\begin{cousinweb}
\textbf{\ur{درہ}} (\emph{darrā}) is usually traced to Persian \emph{dara}, \textquotedblleft{}a valley, a mountain pass\textquotedblright{}.
\end{cousinweb}

\begin{grammarlens}[title={What you've built}]
One more word for getting around.
\end{grammarlens}

\subsection*{Your turn}
\begin{itemize}
\raggedright
  \item \emph{Say it:} \emph{darrā}
  \item \emph{Say it:} \emph{darrā}, once more
  \item \emph{Say it:} \emph{sardiyoṅ meṅ darrā barf se band hotā hai.} — In winter the pass is closed by snow.
  \item \emph{Recall:} say the Urdu for a distance, then the Urdu for a kilometre, then say \emph{darrā} again
\end{itemize}

\begin{tcolorbox}[breakable,colback=teal!4,colframe=teal!35!black,title={Before you move on}]
What does \ur{درہ} mean? (\textquotedblleft{}A mountain pass\textquotedblright{}.) Say it once more.
\end{tcolorbox}
\section[\texorpdfstring{\ur{موڑنا} (moṛnā)}{موڑنا (moṛnā)}]{\ur{موڑنا} (moṛnā) — to turn (something); to bend}

\label{lesson:UR-C397-morna}



\begin{quote}
Before the new one: say the Urdu for a kilometre, then the Urdu for a mountain pass.
\end{quote}

\subsection*{You'll want to know: \ur{موڑنا}}
\textbf{\ur{موڑنا}} — \emph{moṛnā} — \textquotedblleft{}to turn (something); to bend\textquotedblright{}.

\begin{quote}
\ur{ڈرائیور} \ur{نے} \ur{موڑ} \ur{پر} \ur{گاڑی} \ur{بائیں} \ur{موڑی۔}
\end{quote}

\begin{quote}
\emph{ḍrāivar ne moṛ par gāṛī bāeṅ moṛī.}
\end{quote}

\begin{quote}
At the bend the driver turned the car to the left.
\end{quote}

Past, with \ur{نے} (he or she …): \ur{اس} \ur{نے} \ur{سائیکل} \ur{موڑی}, \emph{us ne sāikil moṛī}.

\begin{cousinweb}
\textbf{\ur{موڑنا}} (\emph{moṛnā}) is inherited: CDIAL lists the Hindi cousin \emph{moṛnā}, \textquotedblleft{}twists, bends\textquotedblright{}, under Sanskrit \emph{muṭati} and its variant \emph{mōṭati} (CDIAL 10186). \textbf{\ur{مڑنا}} is its partner, to turn by yourself, and \textbf{\ur{موڑ}} is the bend.
\end{cousinweb}

\begin{grammarlens}[title={What you've built}]
One more word for getting around.
\end{grammarlens}

\subsection*{Your turn}
\begin{itemize}
\raggedright
  \item \emph{Say it:} \emph{moṛnā}
  \item \emph{Say it:} \emph{moṛnā}, once more
  \item \emph{Say it:} \emph{ḍrāivar ne moṛ par gāṛī bāeṅ moṛī.} — At the bend the driver turned the car to the left.
  \item \emph{Recall:} say the Urdu for a kilometre, then the Urdu for a mountain pass, then say \emph{moṛnā} again
\end{itemize}

\begin{tcolorbox}[breakable,colback=teal!4,colframe=teal!35!black,title={Before you move on}]
What does \ur{موڑنا} mean? (\textquotedblleft{}To turn (something); to bend\textquotedblright{}.) Say it once more.
\end{tcolorbox}
\section[\texorpdfstring{\ur{گھمانا} (ghumānā)}{گھمانا (ghumānā)}]{\ur{گھمانا} (ghumānā) — to show (someone) around; to turn round}

\label{lesson:UR-C397-ghumana}



\begin{quote}
Before the new one: say the Urdu for a mountain pass, then the Urdu for to turn (something).
\end{quote}

\subsection*{You'll want to know: \ur{گھمانا}}
\textbf{\ur{گھمانا}} — \emph{ghumānā} — \textquotedblleft{}to show (someone) around; to turn round\textquotedblright{}.

\begin{quote}
\ur{ایک} \ur{دوست} \ur{نے} \ur{ہمیں} \ur{شہر} \ur{گھمایا۔}
\end{quote}

\begin{quote}
\emph{ek dost ne hameṅ shahr ghumāyā.}
\end{quote}

\begin{quote}
A friend showed us round the city.
\end{quote}

Past, with \ur{نے} (he or she …): \ur{اس} \ur{نے} \ur{چابی} \ur{گھمائی}, \emph{us ne chābī ghumāī}.

\begin{cousinweb}
\textbf{\ur{گھمانا}} (\emph{ghumānā}) is built on \textbf{\ur{گھومنا}}, \textquotedblleft{}to go round\textquotedblright{}, which comes from \emph{ghummati}, \textquotedblleft{}revolves\textquotedblright{}, a form CDIAL reconstructs (CDIAL 4485): it is the form for making something, or someone, go round.
\end{cousinweb}

\begin{grammarlens}[title={What you've built}]
One more word for getting around.
\end{grammarlens}

\subsection*{Your turn}
\begin{itemize}
\raggedright
  \item \emph{Say it:} \emph{ghumānā}
  \item \emph{Say it:} \emph{ghumānā}, once more
  \item \emph{Say it:} \emph{ek dost ne hameṅ shahr ghumāyā.} — A friend showed us round the city.
  \item \emph{Recall:} say the Urdu for a mountain pass, then the Urdu for to turn (something), then say \emph{ghumānā} again
\end{itemize}

\begin{tcolorbox}[breakable,colback=teal!4,colframe=teal!35!black,title={Before you move on}]
What does \ur{گھمانا} mean? (\textquotedblleft{}To show (someone) around; to turn round\textquotedblright{}.) Say it once more.
\end{tcolorbox}
